\documentclass[final,5p,times]{elsarticle}
\usepackage{iftex}
\ifPDFTeX
\usepackage[utf8]{inputenc}
\fi
\usepackage[T1]{fontenc}
\usepackage{amsmath}
\usepackage{amssymb}
\usepackage{multirow}
\usepackage{graphicx}
\usepackage{subcaption}
\usepackage{float}
\usepackage{tabularx}
% Manual FloatBarrier definition (placeins package not available)
\makeatletter
\newcommand{\FloatBarrier}{\par\vskip\z@\noindent\mbox{}}
\makeatother
\newfloat{algorithm}{tbp}{loa}
\floatname{algorithm}{Algorithm}
\usepackage{tikz}
\usetikzlibrary{positioning, arrows.meta}
\usepackage{algpseudocode}
\usepackage{hyperref}
\definecolor{ElsevierBlue}{RGB}{0,128,176}
\hypersetup{ colorlinks=true, linkcolor=ElsevierBlue, citecolor=ElsevierBlue, urlcolor=ElsevierBlue
}
\hfuzz=2pt
% ORCID support — renders as a clickable hyperlink in the PDF
\providecommand{\orcid}[1]{\,\href{https://orcid.org/#1}{\textcolor{ElsevierBlue}{\small[ORCID:\,#1]}}}
\hbadness=10000
\setlength{\emergencystretch}{2em}
\makeatletter
\pdfstringdefDisableCommands{% \def\fnref#1{}% \def\corref#1{}% \def\@fnmark{}%
}
\makeatother
\journal{Swarm and Evolutionary Computation} % --- Symbols used in paper ---
\newcommand{\symStatedim}{\mathcal{D}_s}
\newcommand{\symActiondim}{\mathcal{D}_a}
\newcommand{\symPopsize}{N}
\newcommand{\symGamma}{\gamma}
\newcommand{\symMaxIters}{T_{\text{max}}}
\newcommand{\symBufferSize}{C_{\text{rb}}}
\newcommand{\symWindowsize}{W}
\newcommand{\symFeaturecount}{N_f}
\newcommand{\symHiddendimA}{H_a}
\newcommand{\symHiddendimB}{H_b}
\newcommand{\symBatchsize}{B}
\newcommand{\symLearningrate}{\eta}
\newcommand{\symNumepisodes}{E}
\newcommand{\symWarmup}{T_w}
\newcommand{\symNumruns}{N_{\text{runs}}}
\newcommand{\symRhoStag}{\rho_{\text{stag}}}
\newcommand{\symRhoCollapse}{\rho_{\text{coll}}}
\newcommand{\symRhoBonus}{\rho_{\text{div}}}
\newcommand{\symThetaMin}{\theta_{\min}}
\newcommand{\symKappaConv}{\kappa_{c}}
\newcommand{\symEtaCurios}{\eta_{u}}
\newcommand{\symUtd}{R_{\text{utd}}} % --- Experimental values used in paper ---
\newcommand{\valPopsize}{40}
\newcommand{\valMaxIters}{1500}
% Real values from RRSACPSO_Config.m
\newcommand{\valWindow}{5}
\newcommand{\valStateSize}{45}
\newcommand{\valActionSize}{120}
\newcommand{\valBatchSize}{256}
\newcommand{\valBufferSize}{150{,}000}
\newcommand{\valActorLR}{3{\times}10^{-4}}
\newcommand{\valAttnHeads}{4}
\begin{document}
\begin{frontmatter}
\title{Soft Actor‑Critic with Rank-Residual PSO Parameter Adaptation for UAV Path Planning} \author[aff1]{Viet Anh Ngo\orcid{0009-0002-5429-8554}}
\cortext[cor1]{Corresponding author}
\ead{anh.nv215671@sis.hust.edu.vn}
\address[aff1]{School of Mechanical Engineering, Hanoi University of Science and Technology, No.~1 Dai Co Viet Road, Hai Ba Trung District, Hanoi, Vietnam} \begin{abstract}
Particle Swarm Optimization (PSO) has been widely adopted for unmanned aerial vehicle (UAV) path planning due to its conceptual simplicity and fast convergence.
However, its performance is highly sensitive to the choice of hyperparameters, which remain difficult to tune across heterogeneous environments.
This paper presents RRSACPSO, a self-adaptive PSO framework driven by a Soft Actor--Critic controller with rank-residual parameterization that learns to regulate PSO parameters directly from swarm-level feedback.
A compact swarm-statistics state representation summarizes optimization progress, diversity, and stagnation signals without requiring an explicit transformer encoder.
A simple fitness-improvement reward provides direct learning signals tied to changes in global best fitness.
Training stability is supported by twin critics, target networks, and automatic entropy tuning.
The resulting controller emits a low-dimensional latent action that is expanded into particle-wise PSO coefficients through rank-residual control, yielding data-driven adaptation without handcrafted schedules.
Experiments on three-dimensional terrain navigation benchmarks with no-fly zones show that RRSACPSO attains the best mean fitness in two of three RL scenarios, while online RRSACPSO attains the lowest variance across scenarios in our benchmark.
Against traditional self-adaptive PSO baselines, RRSACPSO remains competitive in fitness under a matched online evaluation protocol.
These results indicate that a retained SAC controller with rank-residual PSO adaptation constitutes a stable and practically deployable mechanism for adaptive PSO parameter control in UAV path planning, offering a favorable trade-off between solution quality and computational cost.
\end{abstract} \begin{keyword}
Particle Swarm Optimization \sep Soft Actor-Critic \sep UAV Path Planning \sep Adaptive Parameter Control \sep Swarm Intelligence \sep Rank-Residual Control
\end{keyword}
\end{frontmatter} \section{Introduction}
\label{sec:intro} Unmanned Aerial Vehicles (UAVs) support missions ranging from precision agriculture to defense and disaster response \cite{Delavarpour2021A, Fascista2022Toward, Alotaibi2019LSAR}.
Every mission hinges on path planning to compute safe, constraint-aware trajectories through cluttered three-dimensional environments.
Despite decades of study, planning in large, uncertain airspaces remains challenging because algorithms must simultaneously reason about vehicle dynamics, environmental constraints, and mission-specific trade-offs \cite{AitSaadi2022UAV, Zhou2025A}.
Recent work on coordinated missions highlights the complexity of multiobjective multi-UAV planning, using evolutionary multitasking with adaptive operator selection and knowledge fusion to balance competing mission goals \cite{Meng2025Multiobjective}. Among the rich ecosystem of planners, Particle Swarm Optimization (PSO) \cite{Kennedy1995Particle} has gained widespread adoption for UAV path planning due to its fast convergence and conceptual simplicity \cite{Cheng2024Research}.
However, PSO's performance critically depends on carefully tuned inertia and acceleration coefficients, which are notoriously problem-dependent and can cause premature convergence when left static \cite{Yang2015Lowdiscrepancy}.
While adaptive PSO variants address this through fitness-based or time-varying parameter schedules \cite{Bansal2011Inertia, Liu2025A}, these hand-designed rules struggle to generalize across heterogeneous terrains and obstacle densities. Reinforcement learning (RL) offers a principled alternative by treating parameter adaptation as a sequential decision problem.
Yet existing RL-based PSO controllers exhibit three key limitations: (i) they often rely on instantaneous swarm statistics that ignore temporal dynamics such as stagnation trends and convergence progression, (ii) they usually adapt either a single global coefficient triplet or a discretized parameter set, limiting flexibility, and (iii) they rarely analyze which RL-side design choices materially contribute once the benchmark protocol is matched.
These gaps motivate the present work. We introduce RRSACPSO, a deep reinforcement learning framework that learns PSO control policies directly from swarm-level feedback.
Our contributions are threefold:
\begin{enumerate} \item A retained SAC-based controller with twin critics, automatic entropy tuning, and target networks for stable online parameter adaptation. \item A rank-residual control law that maps a compact latent action into particle-wise PSO coefficients, enabling richer adaptation than a single global coefficient triplet. \item A simple fitness-improvement reward that keeps policy learning aligned with the optimization endpoint throughout training.
\end{enumerate} Extensive experiments on three-dimensional terrain navigation benchmarks with no-fly zones show that RRSACPSO achieves the best mean fitness in two of three RL scenarios, while RRSACPSO yields the lowest run-to-run variance.
The method remains competitive with strong self-adaptive PSO variants under a markedly lower wall-clock cost.
The remainder of this paper is organized as follows.
Section~\ref{sec:background} reviews PSO fundamentals and related adaptive strategies.
Section~\ref{sec:problem} formalizes the path planning problem.
Section~\ref{sec:methods} details the RRSACPSO architecture.
Section~\ref{sec:results} presents experimental results and supplementary training diagnostics.
Section~\ref{sec:discussion} discusses findings and outlines future directions.
\section{Background and Related Work}
\label{sec:background} Path planning for unmanned aerial vehicles (UAVs) has been extensively studied, with methods broadly categorized into classical approaches, sampling-based algorithms, and nature-inspired optimization.
Classical graph-search algorithms such as A* and Dijkstra's algorithm guarantee optimal paths in discretized environments \cite{Hart1968A}.
However, their computational complexity scales poorly with environment resolution, making them impractical for high-dimensional 3D terrain with dynamic constraints.
Potential field methods offer real-time responsiveness but suffer from local minima traps that can strand the UAV in complex obstacle configurations \cite{Mac2016Heuristic}.
Sampling-based planners, including Rapidly-exploring Random Trees (RRT) and Probabilistic Roadmaps (PRM), efficiently explore high-dimensional configuration spaces without explicit discretization \cite{Karur2021A}.
These methods provide probabilistic completeness guarantees but produce suboptimal, jerky paths that require post-processing for smooth UAV trajectories.
Recent hybrid approaches combine sampling with reinforcement learning to improve path quality \cite{Faust2017PRMRL}. Nature-inspired optimization offers an alternative paradigm that directly optimizes path quality metrics.
Genetic Algorithms (GA) evolve populations of candidate paths through selection, crossover, and mutation operators, effectively handling multi-objective constraints such as path length, obstacle clearance, and fuel consumption \cite{Debnath2024An}.
Differential Evolution (DE) and Grey Wolf Optimizer (GWO) have shown competitive performance on UAV routing benchmarks \cite{AitSaadi2022UAV}.
Among swarm intelligence methods, Particle Swarm Optimization (PSO) has emerged as particularly effective for UAV path planning due to its simple implementation, fast convergence, and ability to handle continuous search spaces \cite{Cheng2024Research, Phung2021Safetyenhanced}.
Advanced ensemble strategies have further improved PSO robustness by combining multiple search behaviors within the same swarm \cite{Lynn2017Ensemble}. Particle Swarm Optimization (PSO) is a population-based stochastic search inspired by flocking and swarming behaviors \cite{Kennedy1995Particle}.
A swarm of candidate solutions (particles) moves through the search space by updating velocity and position at each iteration.
Updates combine each particle's own best-found solution (cognitive component) with the swarm's best-found solution (social component), while an inertia term maintains momentum \cite{Clerc2002The}:
\begin{equation}
\left\{
\begin{aligned}
v_i^{t+1} &= w\,v_i^{t} + c_1 r_1^{t}\odot(p_i^{t}-x_i^{t}) + c_2 r_2^{t}\odot(g^{t}-x_i^{t}),\\
x_i^{t+1} &= x_i^{t}+v_i^{t+1}.
\end{aligned}
\right.
\end{equation}
Here, \(w\) is the inertia weight, \(c_1\) and \(c_2\) are the cognitive and social coefficients, and \(r_1,r_2\) are random vectors.
PSO performance critically depends on these control parameters. Parameter adaptation methods fall into three categories: time-varying schedules, fitness-based adaptation, and learning-based control.
Time-varying schedules adjust parameters according to iteration count.
Linearly decreasing inertia weight (LDIW) encourages initial exploration and later exploitation \cite{Eberhart2000Comparing}.
Time-varying acceleration coefficients (TVAC) gradually shift emphasis from cognitive to social learning \cite{Ratnaweera2004Selforganizing}.
Adaptation of inertia weight remains a primary focus, with studies demonstrating that fuzzy-based or oscillating schedules can effectively maintain swarm diversity \cite{Nickabadi2011A}.
Sinusoidal schedules introduce cyclic transitions between exploratory and exploitative phases \cite{Kentzoglanakis2009Particle, Chen2018A}.
While intuitive and easy to implement, these schedules cannot react to swarm-specific dynamics such as stagnation or premature convergence. Fitness-based adaptation uses swarm feedback to dynamically adjust parameters.
Liu et al.\ \cite{Liu2025A} proposed a spherical vector-based adaptive evolutionary PSO (hereafter referred to as SAEPSO) with per-particle adaptation where a fitness adaptation parameter $fap_i^{t}$ normalizes each particle's fitness relative to global and average performance, combined with sinusoidal temporal modulation.
This approach achieves per-particle heterogeneity---better-performing particles receive different parameters than exploratory ones.
Isiet and Gadala's UAPSO \cite{Isiet2019Selfadapting} derives parameters directly from an evolutionary-state metric $ES_i^{t}$, using it simultaneously as inertia weight and to switch between exploration-favoring and exploitation-favoring coefficient schedules.
On benchmark functions, fitness-based methods consistently outperform static schedules \cite{Li2017Particle}.
Zhan et al.\ \cite{Zhan2009Adaptive} established that adaptively partitioning the swarm into different evolutionary states allows the algorithm to dynamically switch between exploration and exploitation based on real-time convergence metrics.
However, convergence analysis reveals that adapted parameters often fail to reach stable points or satisfy theoretical convergence criteria \cite{Harrison2018Selfadaptive}, and predicting long-term performance from short-term observations remains challenging. Entropy-based strategies monitor swarm distribution disorder to prevent premature convergence \cite{Han2022Improved}.
These methods complement fitness-based approaches by providing diversity-aware signals. \subsection{Reinforcement Learning for PSO Parameter Control} Reinforcement learning (RL) offers a principled framework for adaptive parameter control by treating each PSO iteration as a decision point \cite{Sutton2018Reinforcement, Bai2023Evolutionary}.
The agent observes swarm statistics, outputs parameter values, and receives rewards reflecting optimization progress.
This subsection categorizes existing RL-PSO methods and identifies gaps addressed by the proposed framework. Value-based methods discretize the parameter space for Q-learning \cite{Watkins1992Qlearning}.
Meng et al.\ \cite{Meng2023Multistrategy} proposed MPSORL using tabular Q-learning with non-uniform fitness-based state partitioning, achieving multi-strategy selection across exploration-exploitation phases.
DPQ-PSO \cite{Aoun2024Deep} combines Deep Q-Networks with Hidden Markov Models for evolutionary state classification, enabling discrete parameter selection based on inferred swarm phases.
Li et al.\ \cite{Li2023Reinforcement} introduced RLNDM-PSO, which uses Q-learning to select among neighborhood differential mutation strategies, demonstrating that RL can effectively manage operator selection in hybrid PSO variants.
However, discretization necessitates coarse quantization that limits fine-grained control, particularly problematic in continuous parameter spaces where small adjustments can significantly impact convergence behavior. Policy gradient methods enable continuous parameter adaptation without discretization.
RLAM-PSO \cite{Yin2023RLAMPSO} uses Deep Deterministic Policy Gradient (DDPG) \cite{Lillicrap2016Continuous} to map swarm states directly to continuous $(w, c_1, c_2)$ values, outperforming time-varying baselines on benchmark functions.
Building on this, Zhang et al.\ \cite{Zhang2024DDPGPSO} proposed DDPGPSO, which dynamically generates PSO parameters through a DDPG agent that monitors the particle swarm solution set in real-time, reducing manual configuration effort while improving convergence speed.
However, DDPG-based methods often suffer from overestimation bias and function approximation errors in complex landscapes, issues typically addressed in general RL literature through Twin-Delayed algorithms \cite{Fujimoto2018Addressing}.
SAC-SAPSO \cite{vonEschwege2024Soft} demonstrates that entropy-regularized SAC achieves more stable learning than deterministic policies by encouraging exploration of diverse parameter configurations, addressing the brittleness of DDPG in noisy optimization landscapes.
Recent work on Hybrid Reinforcement Learning PSO (HRLPSO) \cite{Liu2023HRLPSO} applies RL to optimize initial population distribution, chaotic adaptive weights, and adaptive learning factors simultaneously, achieving nonlinear control of acceleration coefficients.
Despite these advances, policy-gradient PSO controllers still exhibit three practical limitations in the UAV setting.
First, they usually encode the swarm state as a single-timestep snapshot, discarding short-horizon temporal cues such as stagnation trends, feasibility shifts, and convergence pressure \cite{Lesort2018State}.
Second, they typically emit either a single global coefficient triplet or a discretized parameter choice, which restricts particle-differentiated control.
Third, several hybrid designs rely on additional shaping terms or operator-selection rewards, making it harder to interpret whether policy learning is aligned with the final fitness objective.
The retained RRSACPSO configuration addresses these gaps through a lightweight temporal summary state, a continuous rank-residual action that expands to particle-wise coefficients, and a direct fitness-improvement reward. Table~\ref{tab:rl_pso_comparison} summarizes the key characteristics of existing RL-PSO methods and highlights how RRSACPSO differs from prior work. \begin{table}[htbp]
\centering
\small
\caption{Comparison of RL-based PSO parameter adaptation methods. RRSACPSO combines a SAC-family controller, rank-residual action expansion, and a direct fitness-improvement reward.}
\label{tab:rl_pso_comparison}
\resizebox{\columnwidth}{!}{%
\begin{tabular}{lccccc}
\hline
Method & RL Algorithm & Action & State Encoder & Reward & Target Net \\
\hline
MPSORL \cite{Meng2023Multistrategy} & Q-Learning & Discrete & MLP & Fitness & --- \\
RLNDM-PSO \cite{Li2023Reinforcement} & Q-Learning & Discrete & MLP & Fitness & --- \\
DPQ-PSO \cite{Aoun2024Deep} & DQN & Discrete & HMM+MLP & Fitness & Yes \\
RLAM-PSO \cite{Yin2023RLAMPSO} & DDPG & Continuous & MLP & Fitness & Yes \\
DDPGPSO \cite{Zhang2024DDPGPSO} & DDPG & Continuous & MLP & Fitness & Yes \\
HRLPSO \cite{Liu2023HRLPSO} & Policy Grad. & Continuous & MLP & Fitness & Yes \\
SAC-SAPSO \cite{vonEschwege2024Soft} & SAC & Continuous & MLP & Fitness & Yes \\
\hline
RRSACPSO (Ours) & SAC & Continuous & 15D temporal summary & Fitness & Yes \\
\hline
\end{tabular}%
}
\end{table} RRSACPSO addresses the retained design goals through a compact swarm-state representation, a twin-critic SAC backbone with target networks, rank-residual parameter expansion, and a direct fitness-improvement reward that keeps the learning objective aligned with the optimization target. \section{Problem Formulation}
\label{sec:problem} We consider the problem of planning a collision-free trajectory for a UAV operating in a 3D environment.
The environment is represented by a digital elevation map (DEM) $\mathcal{E}: \mathbb{R}^2 \to \mathbb{R}$ that returns terrain height for any $(x, y)$ coordinate, augmented with a set of cylindrical no-fly zones $\mathcal{O} = \{(c_j, r_j, h_j)\}_{j=1}^{M}$ representing threat areas.
The UAV must navigate from a start position $s_{\text{start}} \in \mathbb{R}^3$ to a goal position $s_{\text{goal}} \in \mathbb{R}^3$ while minimizing the fitness function defined in Equation~\ref{eq:fitness}. \begin{figure}[!htb]
\centering
\includegraphics[width=\linewidth]{figures/ProblemFormulation.pdf}
\caption{Geometric definition of UAV path planning parameters}
\label{fig:problem_formulation}
\end{figure} A candidate path is encoded as a sequence of $n$ waypoints $P = [p_1, p_2, \ldots, p_n]$ where each $p_i = (x_i, y_i, z_i) \in \mathbb{R}^3$.
The UAV follows linear segments between consecutive waypoints, yielding total path length:
\begin{equation}
L(P) = \sum_{i=1}^{n-1} \|p_{i+1} - p_i\|_2
\label{eq:path_length}
\end{equation} Path quality is evaluated through a weighted fitness function that balances competing objectives:
\begin{equation}
\begin{aligned}
f(P) ={}& \lambda_1 L(P) + \lambda_2 C_{\text{turn}}(P) + \lambda_3 C_{\text{climb}}(P) \\
&+ \lambda_4 C_{\text{collision}}(P) + \lambda_5 C_{\text{alt}}(P)
\end{aligned}
\label{eq:fitness}
\end{equation} The turn penalty $C_{\text{turn}}$ penalizes sharp directional changes that exceed UAV maneuverability:
\begin{equation}
C_{\text{turn}}(P) = \sum_{i=2}^{n-1} \max\left(0, \theta_i - \theta_{\max}\right)^2
\end{equation}
where $\theta_i$ is the turn angle at waypoint $i$ and $\theta_{\max}$ is the maximum allowable turn. The climb penalty $C_{\text{climb}}$ penalizes excessive altitude changes:
\begin{equation}
\begin{aligned}
C_{\text{climb}}(P) = \sum_{i=1}^{n-1} \max\Bigl(0, &\frac{|z_{i+1} - z_i|}{\|p_{i+1} - p_i\|_{xy}} \\
&- \gamma_{\max}\Bigr)^2
\end{aligned}
\end{equation}
where $\gamma_{\max}$ is the maximum climb/descent gradient. The collision penalty $C_{\text{collision}}$ combines terrain and obstacle violations:
\begin{equation}
\begin{aligned}
C_{\text{collision}}(P) = \sum_{i=1}^{n} \Bigl[ &\max(0, \mathcal{E}(x_i, y_i) + h_{\min} - z_i) \\
&+ \sum_{j=1}^{M} \max(0, r_j + d_{\text{safe}} - d_{ij}) \Bigr]
\end{aligned}
\end{equation}
where $d_{ij}$ is the distance from waypoint $i$ to obstacle $j$. The altitude cost $C_{\text{alt}}$ encourages the UAV to maintain low-level flight for terrain masking and energy efficiency:
\begin{equation}
C_{\text{alt}}(P) = \sum_{i=1}^{n} z_i
\end{equation} The weights $\lambda_1, \ldots, \lambda_5$ balance path efficiency against constraint satisfaction.
Lower fitness values indicate better paths. \section{Methods}
\label{sec:methods} The RL-based PSO parameter adaptation problem is formulated as a Markov Decision Process (MDP).
In the proposed algorithm, the MDP is defined by the tuple $(\mathcal{S}, \mathcal{A}, \mathcal{R}, \mathcal{T}, \symGamma)$ as illustrated in Figure~\ref{fig:mdp_components}.
The state space $\mathcal{S}$ encodes compact swarm statistics over a short temporal window.
The action space $\mathcal{A}$ is a 9D latent control vector which is expanded into particle-wise PSO coefficients through the rank-residual mapping.
The reward function $\mathcal{R}$ provides direct feedback based on whether the swarm improves the current best fitness.
The transition rules $\mathcal{T}$ describe deterministic updates computed via the standard PSO velocity and position equations, functioning effectively as the environment transition function during learning.
The discount factor $\symGamma$ determines the temporal preference for cumulative returns. \begin{figure}[!htb] \centering \includegraphics[width=\linewidth]{figures/MDP.pdf} \caption{MDP formulation of RRSACPSO. At each PSO iteration $t$, the SAC agent observes compact swarm statistics $s_t$, outputs a latent action that is expanded into particle-wise PSO parameters, and receives a fitness-improvement reward $r_t$ based on whether the swarm improves the current best path.} \label{fig:mdp_components}
\end{figure} \subsection{Compact Swarm-State Representation}
To address the limitations of static state observations used in prior RL-PSO methods, we use a compact swarm-statistics state representation.
At each iteration, the controller extracts nine normalized swarm-level features: iteration progress, spatial diversity ratio, elite-fitness advantage, global stagnation pressure, feasibility-quality readiness, constraint-pressure intensity, attraction toward the actual global best, velocity alignment toward the global best, and particle-level stagnation pressure.
These features are buffered over a temporal window of eight iterations.
In the retained benchmarked configuration, the deployed encoder is intentionally lightweight: it forms a 15D state by concatenating the current 9D feature vector, the first three coordinates of the temporal mean, and the first three coordinates of the one-step trend.
This keeps the controller online-feasible while still exposing short-horizon temporal context. \subsection{Reward Function} The retained controller uses a direct fitness-improvement reward. Let $\Delta f = f_{\text{gbest}}^{t-1} - f_{\text{gbest}}^{t}$. The reward is defined as:
\begin{equation}
r_t =
\begin{cases}
1, & \text{if } \Delta f > 0 \\
-1, & \text{otherwise}
\end{cases}
\label{eq:reward_total}
\end{equation}
This design keeps the reinforcement signal aligned with the optimization objective while avoiding extra reward-shaping terms. \subsection{Soft Actor-Critic Backbone} RRSACPSO employs Soft Actor-Critic \cite{Haarnoja2018Soft} for learning the parameter adaptation policy.
SAC maximizes the entropy-regularized expected return: \begin{equation}
J(\pi) = \mathbb{E}_{\tau \sim \pi}\left[\sum_{t=0}^{T} \gamma^t \left(r_t + \alpha \mathcal{H}(\pi(\cdot|\mathbf{s}_t))\right)\right]
\label{eq:sac_objective}
\end{equation} where $\alpha$ is the temperature parameter (automatically tuned) and $\mathcal{H}$ is the policy entropy.
The algorithm maintains three components: Actor Network $\pi_\theta(\mathbf{a}|\mathbf{s})$: A Gaussian policy with learnable mean and diagonal covariance. In the retained implementation the actor is a compact MLP operating on the flat 15D state vector. Actions are sampled via reparameterization:
\begin{equation}
\mathbf{a}_t = \tanh(\boldsymbol{\mu}_\theta(\mathbf{s}_t) + \boldsymbol{\sigma}_\theta(\mathbf{s}_t) \odot \boldsymbol{\epsilon}), \quad \boldsymbol{\epsilon} \sim \mathcal{N}(0, \mathbf{I})
\end{equation} Twin Critic Networks $Q_{\phi_1}, Q_{\phi_2}$: Two independent scalar Q-functions are retained together with target networks. The critic loss minimizes the Bellman error using the minimum of both critics:
\begin{equation}
\begin{aligned}
L_Q(\phi_i) = \mathbb{E}_{(\mathbf{s},\mathbf{a},r,\mathbf{s}')} \Bigl[ \bigl(Q_{\phi_i}(\mathbf{s}, \mathbf{a}) - (&r + \symGamma \min_{j=1,2} Q_{\phi_j}(\mathbf{s}', \mathbf{a}') \\
&- \alpha \log \pi_\theta(\mathbf{a}'|\mathbf{s}') )\bigr)^2 \Bigr]
\end{aligned}
\end{equation} Temperature Parameter $\alpha$: Automatically adjusted to match a target entropy $\bar{\mathcal{H}} = -\dim(\mathcal{A}) = -\symActiondim$ via:
\begin{equation}
L_\alpha = \mathbb{E}_{\mathbf{s} \sim \mathcal{D}}\left[-\alpha \left(\log \pi_\theta(\mathbf{a}|\mathbf{s}) + \bar{\mathcal{H}}\right)\right]
\end{equation} In the retained benchmarked configuration, target networks remain enabled, automatic entropy tuning is active, and the critics are standard twin scalar Q-networks; TQC, CrossQ, and critic batch normalization are disabled. \subsection{Algorithm Workflow} The complete integration of the retained SAC agent with the PSO search process is summarized in Algorithm \ref{alg:apexpso}. \begin{algorithm}[H] \caption{RRSACPSO Parameter Adaptation} \label{alg:apexpso} \begin{algorithmic}[1] \State \textbf{Input:} Swarm size $\symPopsize$, Max iterations $\symMaxIters$, Replay buffer $\mathcal{D}$ \State \textbf{Initialize:} Swarm positions $x_i^0$, velocities $v_i^0$, Personal bests $p_i^0$, Global best $g^0$ \State \textbf{Initialize:} Actor $\pi_\theta$, Critics $Q_{\phi_{1,2}}$, Temperature $\alpha$ \For{$t = 1$ to $\symMaxIters$} \State Compute compact swarm statistics $s_t$ \State Sample latent action $\mathbf{a}_t \sim \pi_\theta(\cdot|s_t)$ \State Expand $\mathbf{a}_t$ through rank-residual control to particle-wise PSO coefficients $\{w_{i,t}, c_{1,i,t}, c_{2,i,t}\}_{i=1}^{\symPopsize}$ \For{each particle $i \in \{1 \dots \symPopsize\}$} \State Update velocity $v_i^{t} \leftarrow w_{i,t} v_i^{t-1} + c_{1,i,t} r_1 (p_i^{t-1} - x_i^{t-1}) + c_{2,i,t} r_2 (g^{t-1} - x_i^{t-1})$ \State Update position $x_i^{t} \leftarrow x_i^{t-1} + v_i^{t}$ \EndFor \State Evaluate fitness $f(x_i^t)$ and update $p_i^t$ and $g^t$ \State Compute fitness-improvement reward $r_t$ (Eq. \ref{eq:reward_total}) \State Observe $s_{t+1}$ and store transition $(s_t, \mathbf{a}_t, r_t, s_{t+1})$ in $\mathcal{D}$ \If{training session active} \State Sample minibatch $\mathcal{B} \sim \mathcal{D}$ and update critics, actor, and temperature \EndIf \EndFor \State \Return Optimal path $g^t$ \end{algorithmic}
\end{algorithm} \section{Results}
\label{sec:results} \subsection{Experimental Setup} Path quality is measured by global best fitness (Equation \ref{eq:fitness}), which combines path length, turning penalties, climbing penalties, and collision penalties.
Lower fitness indicates better paths. The main benchmark evaluates the retained RRSACPSO configuration in online rank-residual mode (9D latent action) against RL-based and classical PSO baselines.
These are compared against a broad range of traditional and adaptive PSO variants, including PSO-LDIW, FIPS, CLPSO, SAEPSO, and UAPSO, as well as reinforcement learning-based methods: RLAM-PSO (DDPG) \cite{Yin2023RLAMPSO}, DQN-PSO \cite{Aoun2024Deep}, SAC-SAPSO \cite{vonEschwege2024Soft}, MPSORL, and PPO-PSO.
The repository also contains development-stage component experiments, but the benchmark claims in this paper focus on the retained plain-SAC plus rank-residual configuration described in Section~\ref{sec:methods}.
For each algorithm, we conduct $\symNumruns$ independent runs per scenario and report mean and standard deviation on the same endpoint used for optimization (\texttt{fitnessComponents.totalFitness}).
Inferential analysis uses two-sided paired t-tests as the primary test, with two-sided Wilcoxon signed-rank tests as a robustness check.
Within each scenario, all pairwise p-values are adjusted using Holm--Bonferroni family-wise correction.
For each pair, we additionally report the paired mean difference with a 95\% confidence interval, Cohen's $d_z$, and rank-biserial correlation. Table~\ref{tab:hyperparams} summarises all hyperparameters used in the experiments. \begin{table}[htbp]
\centering
\small
\caption{RRSACPSO hyperparameters. Values extracted directly from \texttt{RRSACPSO\_Config.m}.}
\label{tab:hyperparams}
\begin{tabularx}{\columnwidth}{lXr}
\hline
\textbf{Category} & \textbf{Hyperparameter} & \textbf{Value} \\
\hline
\multirow{4}{*}{Fitness (Eq.~\ref{eq:fitness})} & Path weight $\lambda_1$ (norm.~factor 150) & 0.25 \\ & Turning weight $\lambda_2$ (norm.~factor 3) & 0.25 \\ & Climbing weight $\lambda_3$ (norm.~factor 3) & 0.25 \\ & Altitude weight $\lambda_4$ (norm.~factor 50) & 0.20 \\
\hline
\multirow{1}{*}{RL Reward} & Fitness-improvement signal & $\pm 1$ \\
\hline
\multirow{3}{*}{State Encoder} & Temporal window $W$ & 8 iters \\ & State dimension $\mathcal{D}_s$ & 15 \\ & Feature count $N_f$ & 9 \\
\hline
\multirow{3}{*}{Actor Network} & Hidden layers $H_a$ & $[256, 256]$ \\ & Action dimension $\mathcal{D}_a$ & 9 \\ & Activation & ReLU \\
\hline
\multirow{3}{*}{Critic Network} & Hidden layers $H_b$ & $[256, 256, 128]$ \\ & Number of critics & 2 \\ & Target networks & Enabled \\
\hline
\multirow{5}{*}{Training} & Learning rate $\eta$ (actor/critic/$\alpha$) & $3{\times}10^{-4}$ \\ & Batch size $B$ & 128 \\ & Replay buffer $C_{\text{rb}}$ & 10{,}000 (online mode) \\ & Training episodes $E$ & 1 (online mode) \\ & Discount factor $\gamma$ & 0.99 \\
\hline
\multirow{3}{*}{PSO} & Default population $N$ & 100 \\ & Default iterations per episode & 1000 \\ & Benchmark budget & 100 particles, 1000 iterations, 30 runs \\
\hline
\end{tabularx}
\end{table} For reproducibility, we used the following protocol:
\begin{enumerate} \item \textbf{Evaluation environment:} All comparisons use a fixed map size $[100,100,100]$, start point $[10,95,10]$, goal point $[97,2,10]$, and the three \texttt{ChrismasTerrain2} benchmark scenarios with 0, 5, and 10 danger zones. Terrain is loaded from \texttt{data/ChrismasTerrain2.tif}, rescaled to a maximum elevation of 40, and interpolated to a $100\times100$ grid. \item \textbf{Deterministic paired-seed policy:} For run $r$ and scenario $s$, a common base seed $1000r + 100s$ is used across all algorithms, enabling paired statistical comparison under matched stochastic conditions. \item \textbf{Benchmark budget:} All algorithms are evaluated through a single direct global-planning call per run (no replanning), using a benchmark budget of 100 particles, 1000 iterations, 30 runs per scenario, and a 14-worker parallel pool. \item \textbf{Normalized baseline configuration:} Baselines are compared under the same evaluation protocol while preserving each method's core update equations and action/state design. Hyperparameters are fixed before the final evaluation and are not retuned per scenario. \item \textbf{RRSACPSO configuration:} Experiments use online rank-residual mode with a 9D latent action expanded into particle-wise PSO coefficients. The retained controller enables SAC, automatic entropy tuning, target networks, and twin critics, while TQC, critic batch normalization, CrossQ, cross-scale state, SimBa, prioritized replay, and other research branches are disabled. \item \textbf{Statistical reporting:} Each table reports mean and standard deviation over 30 independent runs per scenario on per-run global fitness. Pairwise inference is based on two-sided paired t-tests (primary) and two-sided Wilcoxon signed-rank tests (robustness), with Holm--Bonferroni correction across all pairwise comparisons within each scenario. We report corrected p-values, 95\% CIs for paired mean differences, and effect sizes ($d_z$, rank-biserial $r$).
\end{enumerate} All results reported in Tables~\ref{tab:rl_comparison} and \ref{tab:others_comparison} follow this fairness protocol. \subsection{Algorithm Comparison Results} 
\begin{figure*}[htbp]
\centering
\begin{subfigure}{0.48\textwidth}
\centering
\includegraphics[width=\linewidth]{figures/results/rl_group_3D.pdf}
\caption{RL Methods 3D Trajectory}
\end{subfigure}
\hfill
\begin{subfigure}{0.48\textwidth}
\centering
\includegraphics[width=\linewidth]{figures/results/other_group_3D.pdf}
\caption{Traditional Methods 3D Trajectory}
\end{subfigure}

\vspace{0.5cm}
\begin{subfigure}{0.48\textwidth}
\centering
\includegraphics[width=\linewidth]{figures/results/rl_group_TopView.pdf}
\caption{RL Methods Top View}
\end{subfigure}
\hfill
\begin{subfigure}{0.48\textwidth}
\centering
\includegraphics[width=\linewidth]{figures/results/other_group_TopView.pdf}
\caption{Traditional Methods Top View}
\end{subfigure}
\caption{Algorithm comparison for UAV path planning in Scenario 1.}
\label{fig:scenario1_comparison}
\end{figure*}
Table~\ref{tab:rl_comparison} and Table~\ref{tab:others_comparison} compare the proposed RRSACPSO against RL-based and traditional PSO variants, respectively.
Among RL-based methods, RRSACPSO achieves the lowest mean fitness in Scenarios 1 and 2 (1531.73 and 1581.86), while MPSORL achieves the best mean in Scenario 3 (1543.13).
RRSACPSO exhibits the most stable dispersion, achieving the lowest standard deviation in all three scenarios (51.25, 87.42, and 69.52).
DQN-PSO degrades sharply as environment complexity increases, exhibiting the highest mean fitness and variance across all scenarios. Among traditional PSO variants (Table~\ref{tab:others_comparison}), SAEPSO achieves the lowest mean fitness in Scenarios 1 and 3 (1526.78, 1616.99), while SAEPSO* is best in Scenario 2 (1618.68).
SAEPSO yields the lowest variability in all three scenarios (std 30.92, 54.13, 38.99).
However, RRSACPSO remains competitively close to SAEPSO while offering a decisive computational advantage: wall-clock runtime requires 11--15 seconds per run versus 72--101 seconds for SAEPSO Full, a 7--8$\times$ speedup.
This places RRSACPSO on a favorable trade-off curve when both solution quality and deployment speed are objectives: near-SAEPSO fitness at a fraction of the runtime cost.
PSO-LDIW remains a competitive classical schedule, while FIPS and CLPSO perform poorly in this domain (mean fitness 1859--2566), as their comprehensive learning strategies disperse particles across the search space without sufficient exploitation pressure. PPO-PSO is included as a baseline and represents our own implementation of a PSO controller with a Proximal Policy Optimization policy network, following the foundational architecture principles of Schulman et al.\ \cite{Schulman2017Proximal} for the actor and critic, adapted for continuous parameter output. \begin{table*}[t]
\centering
\small
\resizebox{\textwidth}{!}{%
\begin{tabular}{lcccccc}
\hline
Scenario & RLAMPSO\cite{Yin2023RLAMPSO} & DQN-PSO\cite{Aoun2024Deep} & SAC-SAPSO \cite{vonEschwege2024Soft} & PPO-PSO & MPSORL & RRSACPSO\\
\hline
Scenario 1& 1646.70 $\pm$ 96.27 & 1906.34 $\pm$ 239.09 & 1737.16 $\pm$ 149.80 & 1614.72 $\pm$ 127.81 & 1678.15 $\pm$ 112.90 & \textbf{1531.73} $\pm$ \textbf{88.63} \\
Scenario 2& 1706.57 $\pm$ 98.18 & 1949.08 $\pm$ 234.77 & 1732.93 $\pm$ 111.42 & 1643.29 $\pm$ 96.88 & 1639.52 $\pm$ \textbf{88.98} & \textbf{1581.86} $\pm$ 111.43 \\
Scenario 3& 1745.02 $\pm$ 102.05 & 2105.80 $\pm$ 241.05 & 1765.01 $\pm$ 123.73 & 1609.60 $\pm$ 187.68 & \textbf{1543.13} $\pm$ \textbf{77.23} & 1636.70 $\pm$ 102.37 \\
\hline
\end{tabular}%
}
\caption{Performance comparison of RL-based PSO parameter adaptation algorithms across three scenarios. Results show Mean $\pm$ Standard Deviation of fitness over 30 runs. Bold marks the best per scenario.}
\label{tab:rl_comparison}
\end{table*} \begin{table*}[t]
\centering
\small
\resizebox{\textwidth}{!}{%
\begin{tabular}{lcccccccc}
\hline
Scenario & RRSACPSO & PSO-Standard & PSO-LDIW & FIPS & CLPSO & SAEPSO \cite{Liu2025A} & SAEPSO* \cite{Liu2025A} & UAPSO \cite{Isiet2019Selfadapting} \\
\hline
Scenario 1 & 1859.1 $\pm$ 281.7 & 1746.2 $\pm$ \textbf{140.5} & \textbf{1684.9} $\pm$ 170.4 & 1737.0 $\pm$ 181.9 \\
\hline
\end{tabular}%
}
\caption{Performance comparison of traditional PSO parameter adaptation algorithms across the benchmark scenarios. Results show Mean $\pm$ Standard Deviation of fitness over 30 runs.}
\label{tab:others_comparison}
\end{table*} The supplementary training diagnostics visualized in Figure~\ref{fig:episode_rewards} and Figure~\ref{fig:fitness_convergence} indicate that the retained binary reward provides a usable online learning signal for the SAC controller family, but the primary claims of this paper are based on the matched final benchmark in Tables~\ref{tab:rl_comparison} and \ref{tab:others_comparison}. \subsection{Results Summary} The matched benchmark across RL-based and classical baselines reveals several key findings for adaptive PSO-based UAV path planning: RL Framework Selection. RRSACPSO achieves the lowest mean fitness in two of three scenarios among RL-based controllers, while MPSORL is best in Scenario 3.
Across runs, RRSACPSO has the lowest standard deviation in all three scenarios.
RLAMPSO remains competitive, but DQN-PSO degrades sharply with scenario complexity and shows the largest variance.
The maximum-entropy SAC objective together with the rank-residual parameterization appears to favor stability in the parameter adaptation space. Comparison with Traditional Methods. SAEPSO achieves the lowest mean fitness in Scenarios 1 and 3, while SAEPSO* achieves the best mean in Scenario 2.
SAEPSO also yields the lowest variability, but RRSACPSO has substantially lower wall-clock runtime in the current benchmark.
Classical time-varying schedules (PSO-LDIW) remain strong baselines that should not be overlooked. Practitioner Recommendations. For deployment, RRSACPSO provides strong performance when online adaptation overhead is acceptable, while SAEPSO is preferred when the best mean fitness is the priority and compute budget allows.
PSO-LDIW remains a strong alternative if RL training is not feasible. \section{Discussion}
\label{sec:discussion} The experimental results demonstrate that RRSACPSO effectively addresses the parameter adaptation challenge in PSO for complex UAV path planning.
This section interprets key findings, discusses limitations, and situates our contributions within the broader literature. \subsection{Interpretation of Results} The benchmark evidence supports the retained rank-residual controller as the main distinguishing mechanism of the current method.
Compared to prior RL-PSO methods, RRSACPSO combines a lightweight temporal summary state with particle-wise coefficient expansion, preserving richer adaptation than approaches that emit only a single global coefficient triplet.
The current repository also contains exploratory critic-side variants, but these are not part of the benchmarked default method and are therefore not claimed here as validated contributions.
Unlike RLAM-PSO \cite{Yin2023RLAMPSO} which relies on DDPG's deterministic policies, our entropy-regularized approach maintains exploration throughout training.
DQN-PSO's performance degradation in complex scenarios confirms that discretized parameter spaces fundamentally limit adaptation granularity. \subsection{Limitations} Several limitations warrant acknowledgment.
First, the retained online controller still requires hyperparameter tuning and introduces additional compute during each run, which may not be feasible for rapidly changing mission profiles.
Second, our evaluation focused on static terrain maps; dynamic obstacles and time-varying threat zones remain untested.
Third, while the policy generalizes across the three benchmark scenarios, transfer to significantly different terrain types (e.g., urban canyons, open water) was not evaluated.
Fourth, the fitness function weights ($\lambda_1, \ldots, \lambda_5$) were hand-tuned; automatic weight adaptation could further improve performance.
Fifth, the retained binary reward design is intentionally simple; alternative reward-shaping strategies were not explored in the current repository.
Finally, the retained SAC-style controller still introduces non-trivial online training overhead within each run, which may pose challenges for embedded UAV systems with limited compute. \subsection{Practical Implications} For practitioners, RRSACPSO offers a strong balance of path quality and computational efficiency when online adaptation cost is acceptable.
For scenarios where RL training is infeasible, SAEPSO and PSO-LDIW remain strong alternatives.
The modular architecture also makes it possible to test alternative critic or state-encoding choices without changing the overall PSO evaluation pipeline. \section{Conclusion}
\label{sec:conclusion} This paper presented RRSACPSO, a deep reinforcement learning framework for adaptive PSO parameter control in UAV path planning.
The proposed method integrates three key retained elements: (i) a SAC controller with twin critics, automatic entropy tuning, and target networks, (ii) a rank-residual action mapping that expands a compact latent action into particle-wise PSO coefficients, and (iii) a direct fitness-improvement reward that keeps the learning signal aligned with the optimization endpoint. Extensive experiments on 3D terrain navigation benchmarks with no-fly zones demonstrated that RRSACPSO achieves the best mean fitness in two of three RL scenarios, while RRSACPSO provides the lowest run-to-run variance in our RL benchmark.
Against traditional baselines, RRSACPSO matches or approaches strong self-adaptive PSO variants while running at substantially lower wall-clock cost than SAEPSO Full in the current benchmark.
This efficiency advantage is practically significant even though the reported configuration still learns online within each run, making the comparison conservative with respect to any future frozen-policy deployment mode.
Future work should systematically investigate hyperparameter sensitivity for the retained controller, extend the state representation to natively handle dynamic obstacles, and explore ways to reduce online adaptation cost across diverse operational domains. \section*{Declaration of Competing Interests}
The author declares that he has no known competing financial interests or personal relationships that could have appeared to influence the work reported in this paper. \section*{CRediT Author Contribution Statement}
\textbf{Viet Anh Ngo}: Conceptualization, Methodology, Software, Formal analysis, Investigation, Writing -- original draft, Writing -- review \& editing, Visualization. \section*{Data and Code Availability}
The source code for RRSACPSO, including all algorithm implementations, benchmark configurations, and analysis scripts, is available in the supplementary repository. model weights and raw experimental results are available from the corresponding author upon reasonable request. \begin{figure}[htbp]
\centering
\includegraphics[width=\columnwidth]{figures/training/Training_FitnessConvergence_Comparison.pdf}
\caption{Supplementary training fitness diagnostics for RRSACPSO-family controller variants. These traces illustrate learning behavior during development and are not the primary basis of the final matched benchmark claims.}
\label{fig:fitness_convergence}
\end{figure} \begin{figure}[htbp]
\centering
\includegraphics[width=\columnwidth]{figures/training/Training_EpisodeRewards_Comparison.pdf}
\caption{Supplementary episode-reward diagnostics for RRSACPSO-family controller variants under the retained fitness-improvement reward design.}
\label{fig:episode_rewards}
\end{figure}
\begin{thebibliography}{99}
\bibitem{Delavarpour2021A}
Delavarpour, N., Koparan, C., Nowatzki, J., Bajwa, S., Sun, X.
(2021).
``\href{https://doi.org/10.3390/rs13061204}{A technical study on UAV characteristics for precision agriculture applications and associated practical challenges,}''
\textit{Remote Sensing}, 13(6), 1204. \bibitem{Fascista2022Toward}
Fascista, A., Coluccia, A., Ravazzi, C., Magli, E.
(2022).
``\href{https://doi.org/10.3390/s22051824}{Toward integrated large-scale environmental monitoring using WSN/UAV/crowdsensing: A review of applications, signal processing, and future perspectives,}''
\textit{Sensors}, 22(5), 1824. \bibitem{Alotaibi2019LSAR}
Alotaibi, E.~T., Alqefari, S.~S., Koubaa, A.
(2019).
``\href{https://doi.org/10.1109/access.2019.2912306}{LSAR: Multi-UAV collaboration for search and rescue missions,}''
\textit{IEEE Access}, 7, 55817--55832. \bibitem{AitSaadi2022UAV}
Ait Saadi, A.; Soukane, A.; Meraihi, Y.; Benmessaoud, G.; Mirjalili, S.
``\href{https://doi.org/10.1007/s11831-022-09742-7}{UAV Path Planning Using Optimization Approaches: A Survey,}'' \textit{Archives of Computational Methods in Engineering}, vol. 29, pp. 4233--4284, 2022. \bibitem{Zhou2025A}
Zhou, Y.; Yan, L.; Han, Y.; Xie, H.; Zhao, Y.
``\href{https://doi.org/10.3390/drones9030194}{A Survey on the Key Technologies of UAV Motion Planning,}'' \textit{Drones}, vol. 9, no. 3, p. 194, 2025. \bibitem{Meng2025Multiobjective}
Meng, K., Wu, B., Xin, B., Deng, F., Chen, C.
(2025).
``\href{https://doi.org/10.1016/j.swevo.2025.102145}{Multiobjective multi-UAV path planning via evolutionary multitasking optimization with adaptive operator selection and knowledge fusion,}''
\textit{Swarm and Evolutionary Computation}, 99, 102145. \bibitem{Kennedy1995Particle}
Kennedy, J.; Eberhart, R. \href{https://doi.org/10.1109/ICNN.1995.488968}{Particle swarm optimization}. In Proceedings of the IEEE International Conference Neural Networks, Perth, Australia, 27 November--1 December 1995; pp. 1942--1948. \bibitem{Cheng2024Research}
Cheng, Q.; Zhang, Z.; Du, Y.; Li, Y.
``\href{https://doi.org/10.3390/drones8120701}{Research on Particle Swarm Optimization-Based UAV Path Planning Technology in Urban Airspace,}''
\textit{Drones}, vol. 8, no. 12, p. 701, 2024. \bibitem{Yang2015Lowdiscrepancy}
C.~Yang, W.~Gao, N.~Liu, and C.~Song,
``\href{https://doi.org/10.1016/j.asoc.2015.01.004}{Low-discrepancy sequence initialized particle swarm optimization algorithm with high-order nonlinear time-varying inertia weight,}''
\textit{Applied Soft Computing}, vol.~29, pp.~386--394, 2015. \bibitem{Bansal2011Inertia}
Bansal, J.~C., et al. (2011).
``\href{https://doi.org/10.1109/nabic.2011.6089659}{Inertia weight strategies in particle swarm optimization,}''
\textit{Proceedings of the 2011 Third World Congress on Nature and Biologically Inspired Computing}, pp. 633--640. \bibitem{Liu2025A}
Liu, Y., Zhang, H., Zheng, H., et al. (2025).
``\href{https://doi.org/10.1038/s41598-025-85912-4}{A spherical vector-based adaptive evolutionary particle swarm optimization for UAV path planning under threat conditions,}''
\textit{Scientific Reports}, 15, 2116. \bibitem{Hart1968A}
Hart, P.~E., Nilsson, N.~J., Raphael, B.
(1968).
``\href{https://doi.org/10.1109/tssc.1968.300136}{A formal basis for the heuristic determination of minimum cost paths,}''
\textit{IEEE Transactions on Systems Science and Cybernetics}, 4(2), 100--107. \bibitem{Mac2016Heuristic}
Mac, T.~T., Copot, C., Tran, D.~T., De Keyser, R.
(2016).
``\href{https://doi.org/10.1016/j.robot.2016.08.001}{Heuristic approaches in robot path planning: A survey,}''
\textit{Robotics and Autonomous Systems}, 86, 13--28. \bibitem{Karur2021A}
Karur, K., Sharma, N., Dharmatti, C., Siegel, J.~E.
(2021).
``\href{https://doi.org/10.3390/vehicles3030027}{A survey of path planning algorithms for mobile robots,}''
\textit{Vehicles}, 3(3), 448--468. \bibitem{Faust2017PRMRL}
Faust, A., Ramirez, O., Fiser, M., Oslund, K., Francis, A., Davidson, J., Tapia, L.
(2018).
``\href{https://doi.org/10.1109/ICRA.2018.8461096}{PRM-RL: Long-range robotic navigation tasks by combining reinforcement learning and sampling-based planning,}''
\textit{Proceedings of the IEEE International Conference on Robotics and Automation (ICRA)}, pp.~5113--5120. \bibitem{Debnath2024An}
Debnath, D.; Vanegas, F.; Boiteau, S.; Gonzalez, F.
"\href{https://doi.org/10.3390/drones8070302}{An Integrated Geometric Obstacle Avoidance and Genetic Algorithm TSP Model for UAV Path Planning,}" \textit{Drones}, vol. 8, no. 7, p. 302, 2024. \bibitem{Phung2021Safetyenhanced}
Phung, M.~D., Ha, Q.~P.
(2021).
``\href{https://doi.org/10.1016/j.asoc.2021.107376}{Safety-enhanced UAV Path Planning with Spherical Vector-based Particle Swarm Optimization,}''
\textit{Applied Soft Computing}, 109, 107376. \bibitem{Lynn2017Ensemble}
Lynn, N., Suganthan, P.~N.
(2017).
``\href{https://doi.org/10.1016/j.asoc.2017.02.007}{Ensemble Particle Swarm Optimizer,}''
\textit{Applied Soft Computing}, 55, 533--548. \bibitem{Clerc2002The}
Clerc, M., Kennedy, J.
(2002).
``\href{https://doi.org/10.1109/4235.985692}{The particle swarm---explosion, stability, and convergence in a multidimensional complex space,}''
\textit{IEEE Transactions on Evolutionary Computation}, 6(1), 58--73. \bibitem{Eberhart2000Comparing}
Eberhart, R.~C., Shi, Y.
(2000).
``\href{https://doi.org/10.1109/CEC.2000.870279}{Comparing inertia weights and constriction factors in particle swarm optimization,}''
\textit{Proceedings of the IEEE Congress on Evolutionary Computation}, vol. 1, pp. 84--88. \bibitem{Ratnaweera2004Selforganizing}
Ratnaweera, A., Halgamuge, S.~K., Watson, H.~C.
(2004).
``\href{https://doi.org/10.1109/tevc.2004.826071}{Self-organizing hierarchical particle swarm optimizer with time-varying acceleration coefficients,}''
\textit{IEEE Transactions on Evolutionary Computation}, 8(3), 240--255. \bibitem{Nickabadi2011A}
Nickabadi, A., Ebadzadeh, M.~M., Safabakhsh, R.
(2011).
``\href{https://doi.org/10.1016/j.asoc.2011.01.037}{A novel particle swarm optimization algorithm with adaptive inertia weight,}''
\textit{Applied Soft Computing}, 11(8), 3658--3670. \bibitem{Kentzoglanakis2009Particle}
Kentzoglanakis, K., Poole, M.
(2009).
``\href{https://doi.org/10.1145/1569901.1570140}{Particle swarm optimization with an oscillating inertia weight,}''
In: \emph{Proceedings of GECCO '09}, pp. 1749--1750, ACM. \bibitem{Chen2018A}
Chen, K., Zhou, F., Yin, L., Wang, S., Wang, Y., Wan, F.
(2018).
``\href{https://doi.org/10.1016/j.ins.2017.09.015}{A hybrid particle swarm optimizer with sine cosine acceleration coefficients,}''
\textit{Information Sciences}, 422, 218--241. \bibitem{Isiet2019Selfadapting}
Isiet, M., Gadala, M.
(2019).
``\href{https://doi.org/10.1016/j.asoc.2019.105653}{Self-adapting control parameters in particle swarm optimization,}''
\textit{Applied Soft Computing}, 83, 105653. \bibitem{Li2017Particle}
Li, S.-F., Cheng, C.-Y.
(2017).
``\href{https://doi.org/10.1016/j.cie.2017.06.006}{Particle swarm optimization with fitness adjustment parameters,}''
\textit{Computers \& Industrial Engineering}, 113, 831--841. \bibitem{Zhan2009Adaptive}
Zhan, Z.-H., Zhang, J., Li, Y., Chung, H.~S.-H.
(2009).
``\href{https://doi.org/10.1109/tsmcb.2009.2015956}{Adaptive Particle Swarm Optimization,}''
\textit{IEEE Transactions on Systems, Man, and Cybernetics, Part B (Cybernetics)}, 39(6), 1362--1381. \bibitem{Harrison2018Selfadaptive}
Harrison, K.~R., Engelbrecht, A.~P., and Ombuki-Berman, B.~M.
(2018).
\href{https://doi.org/10.1007/s11721-017-0150-9}{\emph{Self-adaptive particle swarm optimization: a review and analysis of convergence}},
Swarm Intelligence,
vol. 12, no. 3, pp. 187--226, \bibitem{Han2022Improved}
Han, Z., Zhang, J., Zhai, Y., Lu, J.
(2022).
``\href{https://doi.org/10.3390/e24010048}{Improved Particle Swarm Optimization Based on Entropy and Its Application in Implicit Generalized Predictive Control,}''
\textit{Entropy}, 24(1), 46. \bibitem{Sutton2018Reinforcement}
Sutton, R.~S., Barto, A.~G.
(2018).
``\href{http://incompleteideas.net/book/the-book-2nd.html}{Reinforcement Learning: An Introduction,}''
\textit{MIT Press}, 2nd ed. \bibitem{Bai2023Evolutionary}
Bai, H., Cheng, R., Jin, Y.
(2023).
``\href{https://arxiv.org/abs/2108.11863}{Evolutionary reinforcement learning: A survey,}''
\textit{arXiv preprint arXiv:2108.11863}. \bibitem{Watkins1992Qlearning}
Watkins, C.~J.~C.~H., Dayan, P.
(1992).
``\href{https://doi.org/10.1007/bf00992698}{Q-learning,}''
\textit{Machine Learning}, 8(3--4), 279--292. \bibitem{Meng2023Multistrategy}
Meng, X., Li, H., and Chen, A.
(2023).
\href{https://doi.org/10.3934/mbe.2023373}{\emph{Multi-strategy self-learning particle swarm optimization algorithm based on reinforcement learning}},
Mathematical Biosciences and Engineering,
vol. 20, no. 5, pp. 8498--8530, \bibitem{Aoun2024Deep}
Aoun, O.
(2024).
``\href{https://doi.org/10.3390/modelling5040089}{Deep Q-Network-Enhanced Self-Tuning Control of Particle Swarm Optimization,}''
\textit{Modelling}, 5(4), 1709-1728, \bibitem{Li2023Reinforcement}
Li, W., Liang, P., Sun, B., Sun, Y., Huang, Y.
(2023).
``\href{https://doi.org/10.1016/j.swevo.2023.101274}{Reinforcement learning-based particle swarm optimization with neighborhood differential mutation strategy,}''
\textit{Swarm and Evolutionary Computation}, 78, 101295. \bibitem{Yin2023RLAMPSO}
Yin, J., Wang, N.
(2023).
``\href{https://doi.org/10.1007/s40747-023-01012-8}{RLAM-PSO: A reinforcement learning-based adaptive matching particle swarm optimization,}''
\textit{Complex \& Intelligent Systems}, 9, 5305--5319. \bibitem{Lillicrap2016Continuous}
Lillicrap, T.~P., Hunt, J.~J., Pritzel, A., Heess, N., Erez, T., Tassa, Y., Silver, D., Wierstra, D.
(2016).
``\href{https://arxiv.org/abs/1509.02971}{Continuous control with deep reinforcement learning,}''
\textit{International Conference on Learning Representations (ICLR)}. \bibitem{Zhang2024DDPGPSO}
Zhang, Y., Liu, H., Wang, X.
(2024).
``\href{https://doi.org/10.1016/j.swevo.2024.101458}{Particle swarm optimization algorithm based on deep deterministic policy gradient,}''
\textit{Swarm and Evolutionary Computation}, 85, 101458. \bibitem{Fujimoto2018Addressing}
Fujimoto, S., van Hoof, H., Meger, D.
(2018).
``\href{https://arxiv.org/abs/1802.09477}{Addressing function approximation error in actor-critic methods,}''
\textit{Proceedings of the 35th International Conference on Machine Learning}, pp. 1587--1596. \bibitem{vonEschwege2024Soft}
von Eschwege, D., and Engelbrecht, A.
(2024).
\href{https://doi.org/10.3390/math12223481}{\emph{Soft Actor-Critic Approach to Self-Adaptive Particle Swarm Optimisation}},
Mathematics,
vol. 12, no. 22, pp. 3481, \bibitem{Liu2023HRLPSO}
Liu, Y., Zhang, J., Chen, W.
(2023).
``\href{https://doi.org/10.3390/math11194089}{Hybrid reinforcement learning particle swarm optimization with chaotic adaptive weight and adaptive learning factors,}''
\textit{Mathematics}, 11(19), 4089. \bibitem{Lange2024Evolution}
Lange, R.~T., Schaul, T., Chen, Y., Zahavy, T., Dalibard, V., Lu, C., Singh, S., Flennerhag, S.
(2024).
``\href{https://arxiv.org/abs/2403.02985}{Evolution Transformer: In-context evolutionary optimization,}''
\textit{arXiv:2403.02985}. \bibitem{Cenikj2024TransOptAS}
Cenikj, G., Lang, R., Engelbrecht, A.~P., Doerr, C., Korosec, P., Eftimov, T.
(2024).
``\href{https://doi.org/10.1145/3638529.3654027}{TransOptAS: Transformer-based algorithm selection for single-objective optimization,}''
\textit{Proceedings of GECCO '24}, pp. 719--727. \bibitem{AlRifaie2017Swarmic}
al-Rifaie, M.~M., Bishop, J.~M.
(2017).
``\href{https://doi.org/10.1007/978-3-319-55849-3_6}{Swarmic sketches and attention mechanisms,}''
In: \emph{Computational Intelligence in Music, Sound, Art and Design}, pp. 85--100, Springer. \bibitem{Chen2021Decision}
Chen, L., Lu, K., Rajeswaran, A., Lee, K., Grover, A., Laskin, M., Abbeel, P., Srinivas, A., Mordatch, I.
(2021).
``\href{https://arxiv.org/abs/2106.01345}{Decision Transformer: Reinforcement Learning via Sequence Modeling,}''
\textit{Advances in Neural Information Processing Systems}, 34, 15084--15097. \bibitem{Lesort2018State}
Lesort, T.; D{\'\i}az-Rodr{\'\i}guez, N.; Goudou, J. F.; Filliat, D. \href{https://doi.org/10.1016/j.neunet.2018.07.006}{State representation learning for control: An overview}. Neural Networks 2018, 108, 379-392. \bibitem{Mnih2015Humanlevel}
Mnih, V., Kavukcuoglu, K., Silver, D., Rusu, A.~A., Veness, J., Bellemare, M.~G., Graves, A., Riedmiller, M., Fidjeland, A.~K., Ostrovski, G., et al. (2015).
``\href{https://doi.org/10.1038/nature14236}{Human-level control through deep reinforcement learning,}''
\textit{Nature}, 518(7540), 529--533. \bibitem{Schulman2017Proximal}
Schulman, J., Wolski, F., Dhariwal, P., Radford, A., Klimov, O.
(2017).
``\href{https://arxiv.org/abs/1707.06347}{Proximal policy optimization algorithms,}''
\textit{arXiv preprint arXiv:1707.06347}. \bibitem{Bhatt2024CrossQ}
Bhatt, A., Palenicek, D., Belousov, B., Argus, M., Amiranashvili, A., Brox, T., Peters, J.
(2024).
``\href{https://doi.org/10.48550/arXiv.1902.05605}{CrossQ: Batch normalization in deep reinforcement learning for greater sample efficiency and simplicity,}''
\textit{arXiv:1902.05605}. \bibitem{Haarnoja2018Soft}
Haarnoja, T., Zhou, A., Abbeel, P., Levine, S.
(2018).
``\href{https://proceedings.mlr.press/v80/haarnoja18b.html}{Soft actor-critic: Off-policy maximum entropy deep reinforcement learning with a stochastic actor,}''
\textit{Proceedings of the 35th International Conference on Machine Learning}, pp. 1861--1870. \bibitem{Ioffe2015Batch}
Ioffe, S., Szegedy, C.
(2015).
``\href{https://proceedings.mlr.press/v37/ioffe15.html}{Batch normalization: Accelerating deep network training by reducing internal covariate shift,}''
\textit{Proceedings of the 32nd International Conference on Machine Learning}, pp. 448--456. \bibitem{Vaswani2017Attention}
Vaswani, A., Shazeer, N., Parmar, N., Uszkoreit, J., Jones, L., Gomez, A.~N., Kaiser, L., Polosukhin, I.
(2017).
``\href{https://arxiv.org/abs/1706.03762}{Attention is all you need,}''
\textit{Advances in Neural Information Processing Systems}, 30, 5998--6008. \end{thebibliography} \end{document}
